\documentclass[8pt]{extarticle}

% ---------- Layout ----------
\usepackage[a4paper,landscape,margin=0.35cm]{geometry}
\usepackage{multicol}
\usepackage{amsmath,amssymb,mathtools}
\usepackage{enumitem}
\usepackage{xcolor}
\usepackage{tcolorbox}
\usepackage{booktabs}
\usepackage{array}
\usepackage{titlesec}
\usepackage{parskip}

\tcbuselibrary{skins,breakable}

\setlength{\parindent}{0pt}
\setlist[itemize]{itemsep=0pt,topsep=1pt,partopsep=0pt,parsep=0pt,leftmargin=1.1em}
\setlength{\columnsep}{0.5cm}
\setlength{\columnseprule}{0.3pt}

% ---------- Section styling ----------
\titleformat{\section}{\normalfont\bfseries\large\color{blue!55!black}}{}{0pt}{}
\titlespacing{\section}{0pt}{4pt}{2pt}
\titleformat{\subsection}{\normalfont\bfseries\color{red!55!black}}{}{0pt}{}
\titlespacing{\subsection}{0pt}{3pt}{1pt}

% ---------- Box for key formulas ----------
\newtcolorbox{key}{
  colback=blue!4, colframe=blue!45!black, boxrule=0.4pt,
  arc=1pt, left=3pt, right=3pt, top=1.5pt, bottom=1.5pt,
  boxsep=0pt, before skip=2pt, after skip=3pt
}
\newtcolorbox{tab}{
  colback=yellow!6, colframe=black!55, boxrule=0.4pt,
  arc=1pt, left=2pt, right=2pt, top=1.5pt, bottom=1.5pt,
  boxsep=0pt, before skip=2pt, after skip=3pt, breakable
}

\newcommand{\stir}[2]{\genfrac{\{}{\}}{0pt}{}{#1}{#2}}   % Stirling 2nd kind
\newcommand{\stirf}[2]{\genfrac{[}{]}{0pt}{}{#1}{#2}}    % Stirling 1st kind (unsigned)
\newcommand{\falling}[1]{#1^{\underline{\phantom{x}}}}
\newcommand{\note}[1]{{\footnotesize\color{gray!70!black}\emph{#1}}}

\pagestyle{empty}

\begin{document}
\begin{multicols}{4}

\begin{center}
{\Large\bfseries MA2214 Midterm Cheatsheet}\\
{\footnotesize Basic counting $\cdot$ Twelvefold way $\cdot$ Generating functions $\cdot$ Recurrences}
\end{center}

%=================================================================
\section{1. Basic Counting}

\subsection*{(AP) / (MP) -- the two engines}
\begin{key}
\textbf{(AP)} Split into \emph{disjoint} cases, add counts.\\
\textbf{(MP)} Build in \emph{independent ordered steps}, multiply counts.
\end{key}
\note{Intuition: AP = OR (union of disjoint bins), MP = AND (a chain of independent choices).}

\subsection*{(CP) / (DP) -- run them backward}
\begin{itemize}
\item \textbf{(CP)} $|U\setminus A| = |U|-|A|$: count the complement when it's easier.
\item \textbf{(DP)} If a procedure with $m$ outcomes hits every object exactly $r$ times, the number of objects is $m/r$.
\end{itemize}
\note{DP is MP in reverse: overcounting by a constant factor $r$ means dividing by $r$.}

\subsection*{Permutations}
\[ P^n_r = n(n-1)\cdots(n-r+1) = \frac{n!}{(n-r)!} \]
\note{$r$ ordered picks without replacement from $n$ items.}

\subsection*{Circular permutations (unnamed seats)}
\[ Q^n_r = \frac{P^n_r}{r} = \frac{n!}{r(n-r)!}, \qquad Q^n_n=(n-1)! \]
\note{Divide by $r$ rotations (DP). ``Glue a block'': $k$ tied objects $\to$ $k!$ internal orders.}
\begin{itemize}
\item $n$ couples, alternating: $(n-1)!\,n!$
\item $n$ couples, each next to spouse: $2^n(n-1)!$ \ ($n\ge2$; careful at $n=1$: answer is $1$, not $2$)
\end{itemize}

\subsection*{Combinations}
\[ \binom{n}{r} = \frac{P^n_r}{r!} = \frac{n!}{r!(n-r)!} = \binom{n}{n-r} \]
\[\sum_{r=0}^n\binom{n}{r}=2^n \]
\note{Choose $r$ unordered from $n$: build via ordered pick, then DP by $r!$ (kill the order).}

\textbf{Pascal:} $\binom{n}{r}=\binom{n-1}{r-1}+\binom{n-1}{r}$ \note{(split on whether elt.\ $n$ is chosen)}

\textbf{Double-counting identities} (count $2$ ways):
\begin{itemize}
\item $r\binom{n}{r}=n\binom{n-1}{r-1}$
\item $(n-r+1)\binom{n}{r-1}=r\binom{n}{r}$
\item $(n-r)\binom{n}{r}=n\binom{n-1}{r}$
\item $\binom{n}{m}\binom{m}{r}=\binom{n}{r}\binom{n-r}{m-r}$
\end{itemize}
\note{Trick: count pairs (a subset, a flagged element inside/outside it) two different orders.}

\textbf{No two consecutive:} \# of $r$-subsets of $[n]$ with no 2 consecutive $=\binom{n-r+1}{r}$.
\note{Shift $a_i\mapsto a_i-(i-1)$ to kill the gaps -- turns it into an ordinary $r$-subset of a smaller set.}

\textbf{Pairings of $2n$ people} into $n$ unordered pairs:
\[ \frac{(2n)!}{2^n\,n!} = (2n-1)!! \]
\note{Line up, cut into pairs: divide out $2^n$ (order within pair) and $n!$ (order of pairs).}

\subsection*{PIE}
\[ |A\cup B| = |A|+|B|-|A\cap B| \]
\[ |A_1\cup\cdots\cup A_q| = \sum|A_i| - \sum|A_i\cap A_j| \]
\[{}+ \sum|A_i\cap A_j\cap A_k| - \cdots \]
\note{Each element in $k$ sets gets counted $\binom{k}{1}-\binom{k}{2}+\cdots = 1$ time, net.}

\textbf{Euler's $\varphi$:} $n=p_1^{m_1}\cdots p_k^{m_k} \Rightarrow \varphi(n)=n\prod_{i=1}^k\left(1-\dfrac1{p_i}\right)$
\note{PIE on ``divisible by $p_i$''; only distinct primes matter, exponents don't.}

\subsection*{(IP) / (BP)}
\begin{itemize}
\item Injection $A\to B$ $\Rightarrow |A|\le|B|$.
\item Bijection $A\to B$ $\Rightarrow |A|=|B|$.
\end{itemize}
\note{Workhorse: count $A$ by matching it to a set $B$ you already know how to count. Check injective (read object back off image) + surjective (every image hit).}

\textbf{Lattice paths} $X\to Y$ ($m$ E-steps, $n$ N-steps): $\binom{m+n}{n}$ (word in E's/N's).

\textbf{Subsets:} $|\mathcal P(X)|=2^n$ via characteristic vectors in $\{0,1\}^n$.

%=================================================================
\section{2. The Twelvefold Way}
\note{$r$ objects into $n$ boxes; objects/boxes each distinct or identical; no restriction / injective (each box $\le 1$) / surjective (each box $\ge1$).}

\begin{tab}
\centering
\renewcommand{\arraystretch}{1.9}
\begin{tabular}{@{}l c c c@{}}
\toprule
 & No restr. & Inject. & Surject. \\
\midrule
Obj D, Box D & $n^r$ & $P^n_r$ & $n!\stir{r}{n}$ \\
Obj D, Box I & $\sum\limits_{k=0}^n\stir{r}{k}$ & $[r\le n]$ & $\stir{r}{n}$ \\
Obj I, Box D & $\dbinom{r+n-1}{r}$ & $\dbinom{n}{r}$ & $\dbinom{r-1}{n-1}$ \\
Obj I, Box I & $\sum\limits_{k=0}^n p(r,k)$ & $[r\le n]$ & $p(r,n)$ \\
\bottomrule
\end{tabular}
\end{tab}
\note{$[r\le n]$ means $1$ if $r\le n$, else $0$. D=distinct, I=identical.}

\subsection*{Stars and bars (Obj I, Box D)}
\[ \binom{r+n-1}{r} \]
\note{\# arrangements of $r$ identical stars and $(n-1)$ bars.}
\note{Same count answers: nonneg.\ int. solutions of $x_1+\cdots+x_n=r$; multisets of size $r$ from $n$ types ($H^n_r$).}
Surjective version (no box empty): bars go in $r-1$ \emph{internal} gaps only:
\[ \binom{r-1}{n-1} \]

\subsection*{Multisets \& permutations}
\begin{itemize}
\item Permutations of $\{r_1\cdot a_1,\dots,r_n\cdot a_n\}$, $r=\sum r_i$: $\dfrac{r!}{r_1!\cdots r_n!}$
\item $r$-perms.\ of $\{\infty\cdot a_1,\dots,\infty\cdot a_n\}$: $n^r$
\end{itemize}

\subsection*{Stirling numbers, 2nd kind $\stir{r}{n}$}
\emph{\# partitions of $r$ distinct elts.\ into $n$ nonempty unlabelled blocks.}
\begin{key}
$\stir{r}{n}=n\stir{r-1}{n}+\stir{r-1}{n-1}$, \quad $\stir{0}{0}=1$
\end{key}
\note{Split on element $r$: alone in new block ($\stir{r-1}{n-1}$), or joins one of $n$ existing blocks ($n\stir{r-1}{n}$).}
\textbf{Explicit (via PIE on surjections):}
\[ \stir{r}{n} = \frac1{n!}\sum_{j=0}^n(-1)^{n-j}\binom{n}{j}j^r \]
\note{$n!\stir{r}{n}=$ \# surjections $[r]\to[n]$ (label the blocks).}
\textbf{Falling factorial:} $x^{\underline r}=x(x-1)\cdots(x-r+1)$
\[ x^r = \sum_{k=0}^r \stir{r}{k}\, x^{\underline k} \]
\note{Count functions $[r]\to[m]$ two ways: directly $m^r$, or by image size $k$.}
\textbf{Bell number} $B_r=\sum_k\stir{r}{k}$ = \# all partitions of an $r$-set.

\subsection*{Stirling numbers, 1st kind (unsigned) $\stirf{r}{n}$}
\emph{\# permutations of $r$ elts.\ with exactly $n$ disjoint cycles.}
\begin{key}
$\stirf{r}{n}=(r-1)\stirf{r-1}{n}+\stirf{r-1}{n-1}$
\end{key}
\note{Insert elt.\ $r$: start new cycle alone ($\stirf{r-1}{n-1}$), or slot into 1 of $r-1$ arc-positions of an existing cycle ($(r-1)\stirf{r-1}{n}$).}
\textbf{Signed:} $s(r,n)=(-1)^{r-n}\stirf{r}{n}$, and
\[ x^{\underline r} = \sum_{k=0}^r s(r,k)\,x^k \]
\note{2nd-kind matrix $[\stir{r}{k}]$ and 1st-kind matrix $[s(r,k)]$ are inverse to each other.}

\subsection*{Integer partitions $p(r,n)$}
\emph{\# ways to write $r$ as an unordered sum of $n$ positive parts.}
\begin{key}
$p(r,n) = p(r-1,n-1) + p(r-n,n)$
\end{key}
\note{Split on smallest part: $=1$ (delete it, $p(r-1,n-1)$) or $\ge2$ (strip 1 from every part, $p(r-n,n)$).}
\note{Ferrers diagram: rows = parts, left-aligned, decreasing. \textbf{Conjugate} (transpose) swaps \#rows $\leftrightarrow$ largest part $\Rightarrow$ \#partitions of $r$ into $n$ parts $=$ \#partitions of $r$ with largest part $n$.}
No closed form for $p(r,n)$ or $p(r)=\sum_k p(r,k)$ exists (see Ch.\,3 for generating-function attack).

\columnbreak
%=================================================================
\section{3. Binomial Thm \& Generating Functions}

\subsection*{Binomial \& multinomial theorems}
\begin{key}
\[(x+y)^n=\sum_{r=0}^n\binom{n}{r}x^{n-r}y^r\]
\end{key}
\note{Expand unsimplified: each of $n$ factors contributes an $x$ or $y$; a string with $r$ $y$'s is chosen in $\binom nr$ ways.}
Plug in constants to get identities:
\begin{itemize}
\item $x=y=1$: \ $\sum_r\binom{n}{r}=2^n$
\item $x=1,y=-1$: \ $\sum_r(-1)^r\binom{n}{r}=0$ \ ($n\ge1$)
\item $\Rightarrow$ even-index sum $=$ odd-index sum $=2^{n-1}$
\item differentiate \& set $x=1$: $\sum_r r\binom{n}{r}=n\,2^{n-1}$
\end{itemize}
\textbf{Multinomial:}
\[(x_1{+}\cdots{+}x_m)^n=\sum \binom{n}{n_1,\dots,n_m}x_1^{n_1}\cdots x_m^{n_m}\]
\[\binom{n}{n_1,\dots,n_m}=\frac{n!}{n_1!\cdots n_m!}\]
\note{Distribute $n$ distinct positions among $m$ letters; $n_1{+}\cdots{+}n_m=n$.}

\subsection*{Generating functions -- setup}
$F(x)=\sum_n a_nx^n$ encodes sequence $(a_n)$; $[x^r]F(x)=a_r$.
Cauchy product: $FG\leftrightarrow c_r=\sum_{k=0}^r a_kb_{r-k}$.

\begin{key}
\textbf{Product interpretation:} if $x_i\in A^{(i)}$ has factor $F^{(i)}(x)=\sum_{k\in A^{(i)}}x^k$, then
\[[x^r]\,F^{(1)}(x)\cdots F^{(m)}(x)\]
\[= \#\{x_1+\cdots+x_m=r : x_i\in A^{(i)}\}\]
\end{key}
\note{Weighted version: use $x^{ik}$ for a variable with coefficient $i$ (e.g.\ coin values, dice weights).}

\textbf{Encoding dictionary:}
\begin{tab}
\centering\renewcommand{\arraystretch}{1.5}
\begin{tabular}{@{}l c@{}}
\toprule
constraint & factor \\ \midrule
none & $1/(1-x)$ \\
$x_i\ge a$ & $x^a/(1-x)$ \\
$x_i\le b$ & $(1-x^{b+1})/(1-x)$ \\
$a\le x_i\le b$ & $(x^a-x^{b+1})/(1-x)$ \\
$x_i$ even & $1/(1-x^2)$ \\
$x_i$ mult.\ of $k$ & $1/(1-x^k)$ \\
\bottomrule
\end{tabular}
\end{tab}

\subsection*{Extracting coefficients (the toolkit)}
\begin{itemize}
\item \textbf{Cumulative sum:} $\dfrac{1}{1-x}F(x) \leftrightarrow c_r=a_0+\cdots+a_r$
\item \textbf{Stars and bars:} $\dfrac{1}{(1-x)^n}=\displaystyle\sum_{r\ge0}\binom{r+n-1}{n-1}x^r$
\item \textbf{Shifting:} $[x^r]\,x^mF(x)=[x^{r-m}]F(x)$
\item \textbf{Partial fractions, distinct roots:} $\dfrac{P(x)}{\prod(1-r_ix)}=\sum\dfrac{A_i}{1-r_ix}$, and $\dfrac1{1-rx}=\sum r^nx^n$.
\item \textbf{Repeated root, mult.\ $m$:} $\dfrac1{(1-rx)^j}=\sum_n\binom{n+j-1}{j-1}r^nx^n$
\end{itemize}
\note{Recipe: (1) write down $F(x)$ from the dictionary, (2) simplify/factor, (3) apply the toolkit.}

\subsection*{Integer partitions via GF}
\[ \sum_{r\ge0}p(r)x^r = \prod_{k\ge1}\frac{1}{1-x^k} \]
\note{Factor $k$ tracks $m_k=$ \# parts equal to $k$; multiplicities are the exponents.}
\textbf{Euler:} \# partitions into distinct parts $=$ \# partitions into odd parts.
\note{Proof: $D(x)=\prod(1+x^k)=\prod\frac{1-x^{2k}}{1-x^k}$ telescopes to $\prod\frac1{1-x^{2m-1}}=O(x)$.}

\subsection*{Exponential generating functions (EGF)}
\[ \sum_{r\ge0}a_r\frac{x^r}{r!} \]
\note{OGF $\to$ selection/distribution problems (order irrelevant). EGF $\to$ arrangement problems (order matters).}
\begin{key}
\textbf{Product rule:} EGF of ``mix several kinds of objects in an arrangement'' $=$ product of each kind's EGF (binomial convolution: coeff.\ $\sum_i\binom ri a_ib_{r-i}$).
\end{key}
\textbf{Dictionary} (type may appear):
\begin{tab}
\centering\renewcommand{\arraystretch}{1.5}
\begin{tabular}{@{}l c@{}}
\toprule
restriction & EGF factor \\ \midrule
any \# times & $e^x$ \\
at most $n$ times & $1+\frac{x}{1!}+\cdots+\frac{x^n}{n!}$ \\
at least once & $e^x-1$ \\
even \# times & $(e^x+e^{-x})/2$ \\
odd \# times & $(e^x-e^{-x})/2$ \\
\bottomrule
\end{tabular}
\end{tab}
\note{Read off coefficients using $e^{nx}=\sum_r n^r\,x^r/r!$. This is PIE done by algebra.}

\columnbreak
%=================================================================
\section{4. Recurrence Relations}

\subsection*{From counting problem to recurrence}
\note{Find the size parameter $n$; look at the ``last piece'' (last tile, largest disk, last element); split into disjoint cases by (AP); get initial conditions from small cases.}
\textbf{Tiling $1\times n$ with $1\times1,1\times2$:} $a_{n+2}=a_{n+1}+a_n$ (Fibonacci).\\
\textbf{Tower of Hanoi:} $a_n=2a_{n-1}+1,\ a_1=1 \Rightarrow a_n=2^n-1$.

\subsection*{Linear homogeneous recurrences}
$c_0a_n+c_1a_{n-1}+\cdots+c_ra_{n-r}=0$. Try $a_n=x^n\Rightarrow$
\begin{key}
\textbf{Characteristic eqn:} $c_0x^r+c_1x^{r-1}+\cdots+c_r=0$
\end{key}
\textbf{Main Thm.}
\begin{itemize}
\item \textbf{(I) distinct roots} $\alpha_1,\dots,\alpha_r$: \ $a_n=A_1\alpha_1^n+\cdots+A_r\alpha_r^n$
\item \textbf{(II) root $\alpha_i$ w/ multiplicity $m_i$:} contributes $(A_{i1}+A_{i2}n+\cdots+A_{im_i}n^{m_i-1})\alpha_i^n$
\end{itemize}
\note{Then plug in initial conditions to solve for the $A$'s (linear system).}
\note{Complex roots come in conjugate pairs $re^{\pm i\theta}$; combine to real form $2r^n\cos(n\theta+\phi)$, or just keep as $A(re^{i\theta})^n+\bar A(re^{-i\theta})^n$.}

\textbf{Fibonacci} ($a_0=a_1=1$): roots $\frac{1\pm\sqrt5}2$, Binet's formula
$a_n=\frac1{\sqrt5}\!\left[\left(\frac{1+\sqrt5}2\right)^{n+1}\!\!-\left(\frac{1-\sqrt5}2\right)^{n+1}\right]$.

\subsection*{General (non-homogeneous) linear recurrences}
$c_0a_n+\cdots+c_ra_{n-r}=f(n)$.
\begin{key}
General solution $=$ (homog.\ solution $h(n)$) $+$ (one particular solution $g_0(n)$)
\end{key}
\note{Difference of any two solutions solves the homogeneous equation -- so this always works, like ODEs.}
\textbf{Trial forms for $g_0(n)$} (base $a$; if $a$ is a root of mult. $m$, multiply trial by $n^m$):
\begin{tab}
\centering\renewcommand{\arraystretch}{1.5}
\begin{tabular}{@{}l c@{}}
\toprule
$f(n)$ & trial $g_0(n)$ \\ \midrule
$P_k(n)$ (poly., $a{=}1$) & $Q_k(n)$ \\
$c\,a^n$ & $C\,a^n$ \\
$P_k(n)\,a^n$ & $Q_k(n)\,a^n$ \\
$f_1+f_2$ & $g_{0,1}+g_{0,2}$ (superposition) \\
\bottomrule
\end{tabular}
\end{tab}
\note{Substitute trial into recurrence, match coefficients.}

\subsection*{Generating function method (handles everything at once)}
\textbf{Steps:} (1) mult.\ recurrence by $x^n$, sum over valid $n$ $\Rightarrow$ linear eqn.\ for $A(x)$;\ (2) solve for $A(x)$ (rational function);\ (3) partial fractions $\to$ read off $a_n$.
\begin{key}
Bookkeeping: $\displaystyle\sum_{n\ge k}a_{n-k}x^n=x^kA(x)$, \quad $\displaystyle\sum_{n\ge k}a_nx^n=A(x)-\sum_{i<k}a_ix^i$
\end{key}
\textbf{Common OGFs:} $a_n{=}1\!\to\!\frac1{1-x}$; \ $a_n{=}n\!\to\!\frac{x}{(1-x)^2}$; \ $a_n{=}r^n\!\to\!\frac1{1-rx}$.
\note{Denominator of $A(x)$ = characteristic polynomial read backwards ($c_0+c_1x+\cdots+c_rx^r$); a repeated root of mult.\ $m$ shows up as $(1-\alpha x)^m$ in the denominator, matching Main Thm.\ part (II) automatically. Non-homog.\ $f(n)$ just adds to the numerator -- no separate particular-solution guess needed.}

\textbf{Partial fractions recap:}
\[\frac{P(x)}{(1-r_1x)(1-r_2x)}=\frac{A}{1-r_1x}+\frac{B}{1-r_2x}\]
\[A=\frac{P(1/r_1)}{1-r_2/r_1},\qquad B=\frac{P(1/r_2)}{1-r_1/r_2}\]

\subsection*{Systems of recurrences}
\note{When one sequence's recurrence needs an auxiliary sequence, either (a) eliminate the auxiliary variable algebraically to get one higher-order recurrence, or (b) write both as generating-function equations and solve the linear system for $A(x)$ (usually cleaner). Both sequences share the same denominator (hence the same characteristic roots / growth rate); only the numerator (initial data) differs.}

\subsection*{Quick sanity checklist}
\begin{itemize}
\item AP: are the cases really \emph{disjoint} and \emph{exhaustive}?
\item MP: are the choices really \emph{independent} of earlier ones?
\item DP/circular: does every group really have the \emph{same} size $r$? (couples $n{=}1$ trap!)
\item GF: did you get the encoding factor right per variable \emph{before} multiplying?
\item Recurrence: does the trial form need an extra $n^m$ because the base is a root?
\end{itemize}

\end{multicols}
\end{document}